\chaves{Incêndios florestais; Visão computacional; Edge AI; LoRa; Nvidia Jetson Orin Nano; Monitoramento ambiental.} % Palavras chave separadas por ponto e vírgula

% Não deixe linha em branco após \begin{resumo}
\begin{resumo}
A incidência persistente de queimadas representa uma grave ameaça socioambiental e econômica no Brasil, com especial gravidade no bioma Cerrado, que concentrou 68,5\% das áreas queimadas no país em 2025, predominantemente em unidades de conservação. Embora existam soluções tecnológicas para o monitoramento remoto de focos de incêndio baseadas em inteligência artificial (IA), sensoriamento e visão computacional, os sistemas atuais enfrentam sérios gargalos operacionais: limitações de memória em dispositivos de borda (Edge AI), restrições de conectividade em regiões remotas desprovidas de cobertura móvel e, sobretudo, custos elevados de implantação e manutenção, que chegam a superar os R\$ 8 milhões, para aplicações em maior escala, em alternativas comerciais baseadas em torres de monitoramento. Diante desse cenário, este trabalho propõe uma arquitetura de baixo custo e alta eficiência para a detecção precoce de incêndios florestais em Áreas de Preservação Ambiental. A solução desenvolvida integra uma câmera PTZ (Pan-Tilt-Zoom) para varredura e visualização da área monitorada, um computador de borda Nvidia Jetson Orin Nano para a execução local de modelos de visão computacional alimentados por um sistema autônomo fotovoltaico. Para contornar os desafios de comunicação, o sistema emprega uma topologia híbrida de conectividade que combina redes 4G e comunicação via rádio de longo alcance (LoRa). O ecossistema é articulado por um software dedicado ao envio automatizado de alertas e ao direcionamento tático de brigadas de combate.
\end{resumo}
